\documentclass[12pt,a4paper]{report}
\usepackage[a4paper,margin=1in]{geometry}
\usepackage{booktabs}
\usepackage{longtable}
\usepackage{array}
\usepackage{graphicx}
\usepackage{svg}
\usepackage{float}
\usepackage{hyperref}
\usepackage{siunitx}

\title{ChatCatHat\texttrademark: Chapter 5 Results and Evaluation}
\author{Privacy Calling Project}
\date{April 10, 2026}

% If your charts are moved to ../data/matrix/charts, update this macro.
\newcommand{\MatrixChartsDir}{../data/matrix-1775300745300/charts}
\newcommand{\PredictiveChartsDir}{../data/predictive-1776567787891/charts}

\begin{document}
\maketitle
\tableofcontents

\setcounter{chapter}{4}
\chapter{RESULTS AND EVALUATION}
\label{chap:results-evaluation}

\section{Evaluation Objectives and Methodology}
\label{sec:eval-objectives}
This chapter evaluates ChatCatHat\texttrademark\ from three aspects: functional correctness, stress behavior under rising pressure, and bottleneck localization using runtime telemetry.

The evaluation data was generated from the repository benchmark workflow:
\begin{itemize}
  \item \texttt{scripts/benchmark-matrix.js}
  \item \texttt{scripts/benchmark-capacity.js}
  \item \texttt{scripts/generate-matrix-charts.js}
\end{itemize}

The analyzed stress dataset is:
\begin{itemize}
  \item \texttt{data/matrix-1775300745300/matrix-runs.json}
  \item \texttt{data/matrix-1775300745300/matrix-rows.csv}
  \item \texttt{data/matrix-1775300745300/matrix-aggregate.json}
  \item \texttt{data/matrix-1775300745300/charts/*}
\end{itemize}

The matrix benchmark uses:
\begin{itemize}
  \item virtual pool sizes \(V \in \{10,20,30\}\)
  \item online users \(U \in \{20,40,60,80,100\}\)
  \item repeats per point \(R=2\)
  \item calls per user \(=2\), max parallel \(=30\)
  \item call timeout \(=25\) s, poll interval \(=1000\) ms
  \item minimum stage duration \(=8000\) ms, cooldown \(=1500\) ms
\end{itemize}

The benchmark window was generated between 2026-04-04T11:06:32Z and 2026-04-04T11:10:31Z.
Across all points and repeats, the total planned calls were 3600.

The core metrics follow the same definitions used in the scripts:
\begin{equation}
\text{CreateSuccessRate}=\frac{\text{create\_success}}{\text{create\_success}+\text{create\_failed}}
\end{equation}
\begin{equation}
\text{TerminalCompletedRate}=\frac{\text{terminal\_completed}}{\text{terminal\_completed}+\text{terminal\_failed}+\text{terminal\_timeout}}
\end{equation}

\section{Functional Correctness Verification}
\label{sec:functional-correctness}
Automated verification was executed on April 10, 2026 with \texttt{npm test}. The result snapshot is shown in Table~\ref{tab:test-summary}.

\begin{table}[H]
\centering
\caption{Automated test result snapshot.}
\label{tab:test-summary}
\begin{tabular}{lc}
\toprule
Item & Value \\
\midrule
Total test files & 15 \\
Passed & 15 \\
Failed & 0 \\
Cancelled/Skipped/Todo & 0 \\
Total duration & 949.84 ms \\
\bottomrule
\end{tabular}
\end{table}

The passed tests cover API routing, call service behavior, state machine transitions, identity policy, database persistence, capacity model functions, authentication flow, and operations metrics/service management. This confirms that the measured stress behavior is based on a functionally consistent code baseline.

\section{Matrix Stress Test Results}
\label{sec:matrix-results}

\subsection{Overall Outcome}
\label{subsec:overall-outcome}
From the 3600 planned calls:
\begin{itemize}
  \item 617 were created successfully and 2983 failed at create stage
  \item create success ratio over all requests: \(617/3600=0.1714\)
  \item among created calls, terminal results were 577 completed and 40 failed
  \item terminal completed ratio over created calls: \(577/(577+40)=0.9352\)
\end{itemize}

The sampled create failure reasons in \texttt{matrix-runs.json} are consistently \texttt{No available virtual number in pool}, matching the policy behavior in \texttt{src/policy/consistent-callee-policy.js} when no fallback virtual number exists.

\subsection{Create Success and Completion Trends}
\label{subsec:success-trends}
Table~\ref{tab:virtual-summary} summarizes the average trend by virtual pool size and Figure~\ref{fig:create-success-line} to Figure~\ref{fig:marginal-gain} visualize the pressure behavior.

\begin{table}[H]
\centering
\caption{Average results by virtual pool size (from \texttt{charts/summary.md}).}
\label{tab:virtual-summary}
\begin{tabular}{lccc}
\toprule
Virtual Count & Avg Create Success Rate & Avg Terminal Completed Rate & Avg Peak RSS (MB) \\
\midrule
10 & 0.0996 & 0.9689 & 74.00 \\
20 & 0.2277 & 0.9386 & 71.97 \\
30 & 0.3706 & 0.9320 & 64.67 \\
\bottomrule
\end{tabular}
\end{table}

\begin{figure}[H]
\centering
\includesvg[width=0.95\linewidth]{\MatrixChartsDir/01-create-success-rate}
\caption{Create success rate vs online users under different virtual pool sizes.}
\label{fig:create-success-line}
\end{figure}

\begin{figure}[H]
\centering
\includesvg[width=0.95\linewidth]{\MatrixChartsDir/02-terminal-completed-rate}
\caption{Terminal completed rate vs online users under different virtual pool sizes.}
\label{fig:terminal-completed-line}
\end{figure}

\begin{figure}[H]
\centering
\includesvg[width=0.95\linewidth]{\MatrixChartsDir/06-create-success-heatmap}
\caption{Create success heatmap (rows: virtual pool size, columns: online users).}
\label{fig:create-heatmap}
\end{figure}

\begin{figure}[H]
\centering
\includesvg[width=0.95\linewidth]{\MatrixChartsDir/08-marginal-gain-virtual-pool}
\caption{Marginal create-success gain when increasing virtual pool size.}
\label{fig:marginal-gain}
\end{figure}

Key observations are:
\begin{itemize}
  \item Create success rate decreases with user pressure for all virtual pool sizes.
  \item Expanding virtual pool from 10 to 30 significantly improves create success across all user stages.
  \item At 100 users, create success rises from 0.0450 (virtual 10) to 0.1625 (virtual 30), a gain of 0.1175 absolute.
  \item Terminal completed rate remains relatively high (\(\ge 0.8765\)) after call creation, indicating the dominant bottleneck is in front-stage call creation, not post-creation lifecycle completion.
\end{itemize}

\begin{table}[H]
\centering
\caption{Marginal gain in create success rate when expanding virtual pool.}
\label{tab:marginal-gains}
\begin{tabular}{cccc}
\toprule
Online Users & Gain 10\(\rightarrow\)20 & Gain 20\(\rightarrow\)30 & Gain 10\(\rightarrow\)30 \\
\midrule
20 & 0.2875 & 0.3000 & 0.5875 \\
40 & 0.1250 & 0.1375 & 0.2625 \\
60 & 0.0917 & 0.1208 & 0.2125 \\
80 & 0.0813 & 0.0937 & 0.1750 \\
100 & 0.0550 & 0.0625 & 0.1175 \\
\bottomrule
\end{tabular}
\end{table}

Table~\ref{tab:marginal-gains} and Figure~\ref{fig:marginal-gain} show diminishing marginal gain as pressure rises. This suggests that increasing virtual pool size is effective but should be combined with additional scheduling or queueing policies for higher load ranges.

\subsection{Latency Stability}
\label{subsec:latency-stability}
Latency metrics are stable in settle and lifecycle dimensions:
\begin{itemize}
  \item settle P95 range: 2008.5 ms to 2033.0 ms
  \item lifecycle P95 range: 1476.0 ms to 1501.5 ms
\end{itemize}

Create P95 varies more (71.0 ms to 280.0 ms), reflecting increased contention before a call is accepted into the active path.

\begin{table}[H]
\centering
\caption{Latency and terminal quality envelope (matrix aggregate).}
\label{tab:latency-envelope}
\begin{tabular}{lcc}
\toprule
Metric & Minimum & Maximum \\
\midrule
Create P95 (ms) & 71.0 & 280.0 \\
Settle P95 (ms) & 2008.5 & 2033.0 \\
Lifecycle P95 (ms) & 1476.0 & 1501.5 \\
Terminal Completed Rate & 0.8765 & 1.0000 \\
\bottomrule
\end{tabular}
\end{table}

\subsection{Resource Usage and Bottleneck Localization}
\label{subsec:resource-bottleneck}
Resource telemetry indicates the system is not CPU-saturated in this test range:
\begin{itemize}
  \item peak load/core range: 0.0188 to 0.0688
  \item peak RSS range: 57.29 MB to 76.21 MB
  \item peak heap usage ratio range: 0.6118 to 0.9183
\end{itemize}

\begin{figure}[H]
\centering
\includesvg[width=0.95\linewidth]{\MatrixChartsDir/05-server-load-vs-pressure}
\caption{Server load/core vs pressure.}
\label{fig:server-load-line}
\end{figure}

\begin{figure}[H]
\centering
\includesvg[width=0.95\linewidth]{\MatrixChartsDir/03-server-rss-vs-pressure}
\caption{Server peak RSS vs pressure.}
\label{fig:server-rss-line}
\end{figure}

\begin{figure}[H]
\centering
\includesvg[width=0.95\linewidth]{\MatrixChartsDir/09-server-load-heatmap}
\caption{Server load/core heatmap.}
\label{fig:server-load-heatmap}
\end{figure}

\begin{figure}[H]
\centering
\includesvg[width=0.95\linewidth]{\MatrixChartsDir/10-server-rss-heatmap}
\caption{Server RSS heatmap.}
\label{fig:server-rss-heatmap}
\end{figure}

Combined with the sampled failure reason \texttt{No available virtual number in pool}, these measurements show that virtual identity supply and admission policy are the first-order bottlenecks in this dataset, while host CPU and memory remain within safe operational bounds.

\section{AI-Driven Predictive Scaling PoC Evaluation}
\label{sec:predictive-scaling-poc}
To evaluate proactive scheduling, a lightweight Node.js predictive-scaling PoC was introduced:
\begin{itemize}
  \item Holt/EWMA-based traffic forecasting for 15 and 30 minute horizons
  \item Engset-based pool sizing to keep blocking probability under a target threshold
  \item policy comparison under identical synthetic tidal traffic:
  \begin{itemize}
    \item \textbf{Reactive}: scale based on current load only
    \item \textbf{Predictive}: pre-warm pool from short-horizon forecast signal
  \end{itemize}
\end{itemize}

The generated dataset is:
\begin{itemize}
  \item \texttt{data/predictive-1776567787891/series.json}
  \item \texttt{data/predictive-1776567787891/series.csv}
  \item \texttt{data/predictive-1776567787891/summary.json}
  \item \texttt{data/predictive-1776567787891/charts/*}
\end{itemize}

Default PoC settings were: duration \(=360\) min, step \(=1\) min, horizons \(=\{15,30\}\), \(N=300\), hold time \(=1.5\) min, target blocking \(=0.10\), pool bounds \([10,60]\), scale-up delay \(=10\) min, scale-down delay \(=30\) min, and seed \(=42\).

Table~\ref{tab:predictive-summary} summarizes the policy-level KPI.

\begin{table}[H]
\centering
\caption{Reactive vs predictive scaling KPI summary (\texttt{summary.json}).}
\label{tab:predictive-summary}
\begin{tabular}{lcc}
\toprule
Metric & Reactive & Predictive \\
\midrule
Create Success Rate & 0.9156 & 0.9227 \\
Total Create Failures & 542.84 & 497.21 \\
Average Virtual Pool & 30.67 & 30.94 \\
Peak Virtual Pool & 35.00 & 34.00 \\
Blocking Minutes (\(P_b>0.10\)) & 81 & 60 \\
\bottomrule
\end{tabular}
\end{table}

Derived gains from the same run:
\begin{itemize}
  \item failure reduction: \(8.41\%\)
  \item blocking-minute reduction: \(21\) minutes
  \item peak pool did not increase (\(35 \rightarrow 34\))
\end{itemize}

\begin{figure}[H]
\centering
\includesvg[width=0.95\linewidth]{\PredictiveChartsDir/01-traffic-forecast}
\caption{Synthetic tidal traffic with 15/30-minute forecast curves.}
\label{fig:predictive-traffic-forecast}
\end{figure}

\begin{figure}[H]
\centering
\includesvg[width=0.95\linewidth]{\PredictiveChartsDir/02-pool-scaling-compare}
\caption{Virtual pool scaling paths: reactive vs predictive pre-warming.}
\label{fig:predictive-pool-compare}
\end{figure}

\begin{figure}[H]
\centering
\includesvg[width=0.95\linewidth]{\PredictiveChartsDir/03-blocking-compare}
\caption{Blocking probability comparison over time.}
\label{fig:predictive-blocking-compare}
\end{figure}

\begin{figure}[H]
\centering
\includesvg[width=0.95\linewidth]{\PredictiveChartsDir/04-cumulative-failures}
\caption{Cumulative create failures under reactive and predictive policies.}
\label{fig:predictive-cumulative-failures}
\end{figure}

These results indicate that short-horizon predictive pre-warming can reduce create-stage loss with minimal additional average pool usage, supporting the feasibility of an AI-driven proactive scaling path before introducing heavier online ML services.

\subsection{Full Aggregate Matrix}
\label{subsec:full-matrix}
Table~\ref{tab:matrix-aggregate-full} lists the full aggregate matrix for reproducibility.

\begin{table}[H]
\centering
\caption{Matrix stress test aggregate results (2 repeats per point).}
\label{tab:matrix-aggregate-full}
\resizebox{\textwidth}{!}{%
\begin{tabular}{cccccccc}
\toprule
Virtual & Users & Planned Calls & Create Rate & Terminal Completed Rate & Create P95 (ms) & Settle P95 (ms) & Peak RSS (MB)\\
\midrule
10 & 20 & 40 & 0.2250 & 1.0000 & 125.5 & 2012.0 & 76.21\\
10 & 40 & 80 & 0.1000 & 1.0000 & 83.0 & 2011.5 & 72.64\\
10 & 60 & 120 & 0.0750 & 0.9000 & 71.0 & 2014.0 & 73.68\\
10 & 80 & 160 & 0.0531 & 0.9444 & 99.5 & 2011.0 & 73.61\\
10 & 100 & 200 & 0.0450 & 1.0000 & 98.5 & 2008.5 & 73.85\\
20 & 20 & 40 & 0.5125 & 0.9055 & 164.0 & 2021.0 & 69.50\\
20 & 40 & 80 & 0.2250 & 1.0000 & 174.0 & 2027.5 & 71.57\\
20 & 60 & 120 & 0.1667 & 0.9318 & 147.5 & 2019.0 & 72.17\\
20 & 80 & 160 & 0.1344 & 0.9307 & 225.5 & 2028.5 & 72.58\\
20 & 100 & 200 & 0.1000 & 0.9250 & 152.5 & 2021.0 & 74.05\\
30 & 20 & 40 & 0.8125 & 0.9236 & 206.5 & 2033.0 & 74.82\\
30 & 40 & 80 & 0.3625 & 0.9828 & 207.5 & 2030.5 & 57.29\\
30 & 60 & 120 & 0.2875 & 0.9500 & 220.0 & 2028.0 & 61.30\\
30 & 80 & 160 & 0.2281 & 0.8765 & 280.0 & 2022.5 & 63.11\\
30 & 100 & 200 & 0.1625 & 0.9272 & 201.5 & 2024.0 & 66.83\\
\bottomrule
\end{tabular}%
}
\end{table}

\section{Threats to Validity}
\label{sec:threats-validity}
The following limitations should be considered:
\begin{itemize}
  \item The benchmark runs on a single host (\texttt{127.0.0.1}) and does not include WAN jitter/loss.
  \item Each matrix point uses two repeats; higher repeat counts would reduce variance.
  \item Simulated always-online SIP endpoints reduce uncertainty in registration behavior compared to uncontrolled real-user conditions.
  \item Resource measurements are sampled every 2 seconds, so very short spikes may be under-represented.
\end{itemize}

\section{Chapter Summary}
\label{sec:chapter-summary}
The evaluation confirms that ChatCatHat\texttrademark\ is functionally correct on the tested code baseline and that the primary pressure bottleneck is virtual-number admission rather than host CPU saturation.

Under the evaluated matrix:
\begin{itemize}
  \item Larger virtual pools consistently improve create-stage success.
  \item Post-creation terminal completion remains high, indicating stable downstream call lifecycle handling.
  \item Runtime resource levels stay moderate, supporting further scaling work on policy and pool management before infrastructure scaling.
  \item Predictive pre-warming in PoC simulation further improves create-stage success and reduces blocking duration under tidal traffic.
\end{itemize}

\vspace{0.5em}
\noindent
\textbf{Compilation note:} this document uses \texttt{\textbackslash includesvg}. Compile with \texttt{-shell-escape}, or convert the SVG charts in \texttt{\MatrixChartsDir} to PDF and replace with \texttt{\textbackslash includegraphics}.

\end{document}
